% 4-5-2-dm-lago.tex
%  Budget: ~3pp.  HANDOFF 6.5, 6.6, 19.1 (lago_comparison, dm_vs_lago,
%  seed_ensemble).
%  Source map: lago_forecasting_2021 -- published forecasts and paper tables.
%  MANDATE:
%   - headline is "we do NOT beat their best"; never "beats the state of the art"
%   - paper-vs-shipped discrepancy gets its own paragraph (LEAR only)
%   - seed ensemble is SUPPLEMENTARY and must stay visibly separate
%  Rules: decimals in \lr{}; integers Persian; every table cell in \lr{}.

\subsection{مقایسه با نتایج منتشرشده}
\label{sec:4-5-2}

\subsubsection*{چه چیزی با چه چیزی سنجیده می‌شود}

لاگو و همکاران \cite{lago_forecasting_2021} علاوه بر جدول‌های مقاله، فایل‌های
پیش‌بینیِ خام مدل‌هایشان را نیز منتشر کرده‌اند. این پژوهش از همان فایل‌ها استفاده
می‌کند و آن‌ها را با \emph{کد سنجه‌ی خودِ این پژوهش} دوباره امتیازدهی می‌نماید. پس
مقایسه‌ی زیر بر یک بازار، یک پنجره‌ی آزمون، یک مجموعه‌ی مبادی و یک پیاده‌سازیِ واحدِ
سنجه‌ها انجام می‌شود. اگر به‌جای این کار اعداد چاپ‌شده‌ی مقاله در برابر اعداد ما
گذاشته می‌شد، دو مسیرِ محاسبه‌ی متفاوت با هم مقایسه می‌شدند و تفاوت‌های کوچک بی‌معنا
می‌بود.

\begin{table}[htbp]
\centering
\caption{مدل‌های این پژوهش در برابر منتخبی از گونه‌های منتشرشده‌ی لاگو و همکاران
\cite{lago_forecasting_2021}، بر بازار \lr{EPEX-DE} و همان دوره‌ی آزمون. تمام
اعداد با کد سنجه‌ی این پژوهش از فایل‌های پیش‌بینیِ منتشرشده محاسبه شده‌اند.
مرتب‌شده بر پایه‌ی \lr{MAE}.}
\label{tab:lago-comparison}
\small
\setlength{\tabcolsep}{6pt}
\begin{tabular}{llrr}
\toprule
مدل & منبع & \lr{MAE} & \lr{rMAE} \\
\midrule
\lr{DNN Ensemble} & لاگو و همکاران & \lr{\textbf{3.4135}} & \lr{\textbf{0.3740}} \\
ترکیب رژیم‌محور & این پژوهش & \lr{3.5569} & \lr{0.3897} \\
ترکیب ایستا & این پژوهش & \lr{3.5742} & \lr{0.3916} \\
\lr{DNN 4} & لاگو و همکاران & \lr{3.5917} & \lr{0.3935} \\
\lr{LEAR Ensemble} & لاگو و همکاران & \lr{3.6091} & \lr{0.3954} \\
\lr{LSTM} & این پژوهش & \lr{3.8734} & \lr{0.4244} \\
\lr{LEAR-LASSO} & این پژوهش & \lr{3.8988} & \lr{0.4272} \\
\lr{LEAR 1092} & لاگو و همکاران & \lr{3.9298} & \lr{0.4306} \\
\bottomrule
\end{tabular}
\end{table}

\subsubsection*{نتیجه‌ی آزمون، بدون تلطیف}

\cref{tab:dm-lago} چهار مقایسه‌ی مستقیمِ ترکیب رژیم‌محور با گونه‌های لاگو را
می‌آورد.

\begin{table}[htbp]
\centering
\caption{آزمون دیبولد-ماریانو با تصحیح \lr{HAC}، ترکیب رژیم‌محورِ این پژوهش در
برابر گونه‌های منتشرشده. ستون \lr{p} احتمالِ آن است که مدلِ این پژوهش بهتر نباشد.}
\label{tab:dm-lago}
\small
\setlength{\tabcolsep}{6pt}
\begin{tabular}{lrrl}
\toprule
در برابر & \lr{MAE} آن‌ها & \lr{p} & نتیجه \\
\midrule
\lr{DNN Ensemble} & \lr{3.4135} & \lr{0.9873} & آنِ آن‌ها بهتر است \\
\lr{DNN 4} & \lr{3.5917} & \lr{0.3222} & تفاوت معنادار نیست \\
\lr{LEAR Ensemble} & \lr{3.6091} & \lr{0.1271} & تفاوت معنادار نیست \\
\lr{LEAR 1092} & \lr{3.9298} & \lr{1.0e-07} & آنِ ما بهتر است \\
\bottomrule
\end{tabular}
\end{table}

جمله‌ای که باید نوشته شود، همین است و نه بیش از آن: \textbf{این پژوهش بهترین مدل
منتشرشده را شکست نمی‌دهد.} ترکیب \lr{DNN} آن‌ها با \lr{3.4135} در برابر
\lr{3.5569} برتری دارد و این برتری معنادار است (\lr{p} برابر \lr{0.0127} در جهت
مقابل). در برابر دو گونه‌ی دیگرِ آن‌ها \lr{—} \lr{DNN 4} و \lr{LEAR Ensemble}
\lr{—} تفاوت معنادار نیست، یعنی مدل این پژوهش با آن‌ها هم‌تراز است. و بر
\lr{LEAR 1092} برتری روشن دارد.

چنین نتیجه‌ای در فصل نخست پیش‌بینی نشده بود که رخ ندهد؛ برعکس، پس از نقطه‌ی بازبینیِ
هفته‌ی پنجم، دفاع پروژه آگاهانه از مسیر «بهبود جدول امتیازها» به مسیر «مشارکت
روش‌شناختی» منتقل شد. آنچه این بخش نشان می‌دهد \lr{—} هم‌ترازی با میانه‌ی
گونه‌های منتشرشده، با ترکیبی از خانواده‌های مدلیِ متفاوت \lr{—} پشتوانه‌ی همان
تصمیم است، نه شاهدی بر شکست آن.

\subsubsection*{یک یافته‌ی بازتولیدپذیری که ارزش ثبت دارد}

هنگام امتیازدهیِ دوباره‌ی فایل‌های منتشرشده، یک ناهم‌خوانیِ نظام‌مند آشکار شد که
شایسته‌ی ثبت است، چون درباره‌ی خودِ ادبیات مرجع چیزی می‌گوید.

سطرهای \lr{DNN} دقیقاً بازتولید می‌شوند. \lr{DNN Ensemble} در مقاله \lr{3.413}
چاپ شده و از فایل‌ها \lr{3.4135} به دست می‌آید؛ \lr{DNN 4} در مقاله \lr{3.592} و از
فایل‌ها \lr{3.5917}. اختلاف تنها در گِرد‌کردن است.

سطرهای \lr{LEAR} بازتولید نمی‌شوند. \lr{LEAR Ensemble} در جدول مقاله \lr{3.955}
است اما از فایل‌های منتشرشده \lr{3.6091} امتیاز می‌گیرد؛ \lr{LEAR 1092} در مقاله
\lr{4.108} است و از فایل‌ها \lr{3.9298}. الگو یکنواخت است: در تمام گونه‌های
\lr{LEAR}، فایل‌های منتشرشده \emph{بهتر} از جدول چاپ‌شده امتیاز می‌گیرند.

انتخاب این پژوهش، استفاده از فایل‌های منتشرشده بوده است. باید توجه داشت که این
محافظه‌کارانه‌ترین انتخاب ممکن است: فایل‌ها مدل‌های \emph{آن‌ها} را بهتر نشان
می‌دهند، پس فاصله‌ی این پژوهش تا آن‌ها را بیشتر می‌کند، نه کمتر. اگر اعداد چاپ‌شده
مبنا قرار می‌گرفت، ترکیب این پژوهش بر هر دو گونه‌ی \lr{LEAR} برتری آشکار پیدا
می‌کرد. علت این ناهم‌خوانی از بیرون قابل تعیین نیست \lr{—} می‌تواند نسخه‌ای متفاوت
از کد، پنجره‌ای متفاوت یا اشتباهی در نگارش جدول باشد \lr{—} و این پژوهش هیچ‌یک را
ادعا نمی‌کند. آنچه ادعا می‌شود صرفاً همین است که این دو منبع با هم نمی‌خوانند، و
پژوهشگری که به یکی از آن‌ها استناد کند نتیجه‌ای متفاوت از دیگری خواهد گرفت.

\subsubsection*{افزوده‌ی تکمیلی: ترکیب بذرها}

آنچه در ادامه می‌آید \textbf{تکمیلی است و هیچ عددی از آن وارد نتایج اصلی این
پایان‌نامه نمی‌شود.} تمام جدول‌های \cref{sec:4-2} تا \cref{sec:4-5-2}، تحلیل
\cref{sec:4-6} و آزمون برون‌توزیعیِ \cref{sec:5-2} همگی بر مدل \lr{LSTM} با بذر
تثبیت‌شده‌ی ۴۲ استوارند و با این افزوده تغییر نمی‌کنند.

مدل‌های شبکه‌ی عصبی به مقدار اولیه‌ی تصادفی حساس‌اند. برای سنجش اینکه چه مقدار از
فاصله‌ی این پژوهش تا لاگو ناشی از همین حساسیت است، \lr{LSTM} با چهار بذر ۴۲ تا ۴۵
آموزش داده شد و میانگین پیش‌بینی‌های آن‌ها به‌جای مدل تک‌بذری در ترکیب نشست.

\begin{table}[htbp]
\centering
\caption{اثر میانگین‌گیری بر چهار بذرِ \lr{LSTM}. \textbf{تکمیلی} \lr{—} این
اعداد جایگزین نتایج تثبیت‌شده نمی‌شوند و در هیچ جدول دیگری تکرار نمی‌گردند.}
\label{tab:seed-ensemble}
\small
\setlength{\tabcolsep}{6pt}
\begin{tabular}{lrr}
\toprule
پیکربندی & \lr{MAE} & \lr{rMAE} \\
\midrule
\lr{LSTM}، بذر ۴۲ (تثبیت‌شده) & \lr{3.8734} & \lr{0.4244} \\
\lr{LSTM}، ترکیب چهار بذر & \lr{3.6460} & \lr{0.3995} \\
\midrule
ترکیب ایستا، با \lr{LSTM} تثبیت‌شده & \lr{3.5742} & \lr{0.3916} \\
ترکیب ایستا، با ترکیب بذرها & \lr{3.5260} & \lr{0.3863} \\
ترکیب رژیم‌محور، با \lr{LSTM} تثبیت‌شده & \lr{3.5569} & \lr{0.3897} \\
ترکیب رژیم‌محور، با ترکیب بذرها & \lr{3.4994} & \lr{0.3834} \\
\midrule
\lr{DNN Ensemble} (مرجع) & \lr{3.4135} & \lr{0.3740} \\
\bottomrule
\end{tabular}
\end{table}

اثر آن قابل توجه است. خودِ \lr{LSTM} از \lr{3.8734} به \lr{3.6460} می‌رسد، و ترکیب
رژیم‌محور از \lr{3.5569} به \lr{3.4994}. مهم‌تر از عدد، تغییرِ نتیجه‌ی آزمون است:
برتریِ \lr{DNN Ensemble} که با مدل تثبیت‌شده در \lr{p} برابر \lr{0.0127} معنادار
بود، با ترکیب بذرها به \lr{0.0803} می‌رسد و دیگر معنادار نیست.

این نکته را باید با احتیاط خواند. آنچه نشان می‌دهد این نیست که این پژوهش با
\lr{DNN Ensemble} برابر است؛ نشان می‌دهد که بخشی از فاصله‌ی سنجیده‌شده، نوفه‌ی
مقداردهی اولیه است و نه تفاوت روش. ضمناً مقایسه به این شکل کاملاً منصفانه نیست،
چون گونه‌ی مرجع خود یک ترکیب است و مدل تثبیت‌شده‌ی این پژوهش تک‌اجرا. اما دلیل
اینکه این اعداد جای اعداد اصلی را نمی‌گیرند، ساده و انضباطی است: تثبیت نتایج پیش
از آغاز نگارش انجام شد، و جابه‌جاکردنِ پس از آن، هر تضمینی را که آن تثبیت فراهم
می‌کند باطل می‌کرد. کرانِ اوراکل برای پیکربندی ترکیب‌بذری \lr{3.5019} است، در برابر
\lr{3.5580} برای پیکربندی تثبیت‌شده.
